\subsection{Sprint 1}

Na Sprint 1, definimos as \textit{User Stories} e nos dividimos em duas squads. Os alunos puderam escolher uma área de preferência, entre \textit{frontend} e \textit{backend}, e eu escolhi atuar no \textit{frontend}. Essa escolha fazia sentido para mim justamente porque eu tinha pouca experiência nessa parte e queria aproveitar a AGES para aprender algo que ainda não dominava. Sempre tive curiosidade em desenvolvimento de aplicativos, mas nunca tinha tido a oportunidade de trabalhar com React Native, TypeScript ou integração com APIs. Por isso, essa primeira sprint foi um momento de aprendizado intenso e de entender como rodar o \textit{frontend} localmente, como organizar o código e como implementar telas seguindo os protótipos do Figma.

Depois da divisão das squads, os AGES III me ajudaram a fazer o setup do ambiente de desenvolvimento, instalar as dependências e rodar o aplicativo localmente. Eles também me explicaram como funcionava a estrutura do projeto, como os componentes eram organizados. Também me ajudaram a organizar as tarefas e fiquei responsável por desenvolver a tela de acolhimento (onboarding) e criar uma chamada de API para textos. A parte visual da tela foi mais tranquila, pois consegui seguir o protótipo e aplicar os conceitos iniciais de React Native. Já a integração com o \textit{backend} foi mais difícil, porque eu ainda não entendia bem como uma API era chamada a partir do aplicativo e como testar esse fluxo com Postman.

Consegui finalizar a atividade com o apoio de colegas mais experientes, que me ajudaram a entender melhor a comunicação entre \textit{frontend} e \textit{backend}. Essa foi uma das primeiras vezes em que percebi, na prática, como pedir ajuda no momento certo acelera o aprendizado e evita que uma dúvida individual vire bloqueio para a entrega da sprint.
